\chapter{Introduction}
\label{chapter:introduction}

\section{Scientific context and motivation}

Mean-field control (MFC) provides a framework for controlling large populations of interacting agents under a common objective. Such problems arise when the evolution or cost of each agent depends not only on its own state and action, but also on the distribution of the population. For example, in a large computer network, the probability that a machine becomes compromised may depend on the proportion of compromised machines, while a central controller decides how protection resources should be allocated across the network. Directly modelling all agents leads to a state space whose dimension grows with the population size. In the mean-field limit, the empirical distribution is replaced by the law of a representative agent, yielding a control problem whose formulation no longer depends explicitly on the number of agents \cite{motte2022mean}. This reduction makes large systems tractable, but introduces a measure-valued state variable.

Most numerical methods for MFC assume that the transition dynamics and cost functions are known. Depending on the setting, they rely on dynamic programming, forward-backward systems of Hamilton--Jacobi--Bellman and Fokker--Planck equations, or direct differentiation through the population dynamics \cite{pham2018bellman, carmona2018probabilistic, lauriere2021numerical}. In particular, the classical MFC formulation is characterized by a coupled system in which the value function and the population distribution evolve jointly, making the knowledge of the model dynamics central to both the analysis and the numerical solution \cite{pham2018bellman, carmona2018probabilistic}.

Policy-gradient methods are particularly attractive in this setting because they optimize a parametrized policy directly and can accommodate randomized policies, continuous actions, and neural-network parametrizations. In a classical Markov decision process, REINFORCE expresses the gradient of the expected cost using only sampled returns and the score of the policy \cite{williams1992reinforce}. This representation is model-free because, conditionally on the current state and action, the transition kernel does not depend on the policy parameters.

This property no longer holds in mean-field control. The population law is generated by the policy itself: changing the policy parameters changes the population flow, which then changes the policies, costs, and transition kernels evaluated at subsequent times. The classical policy score accounts for the direct effect of the parameters on the sampled actions, but not for their indirect effect through the population law. Consequently, applying REINFORCE without modification generally produces a biased gradient estimator.

A recent method addresses this difficulty in discrete-time, finite-state MFC by perturbing the logits of the population vector and deriving an additional likelihood-ratio correction \cite{meunier2026model}. This establishes that fully model-free policy gradients are possible in this setting. However, the construction relies on one logit coordinate per state and is therefore intrinsically tied to finite state spaces. It does not extend directly to a continuous state space, where the population state is a probability measure with no canonical finite-dimensional parametrization.

The motivation of this work is to develop a perturbation framework that separates the policy-gradient construction from the representation of the population law. In finite state spaces, probability-preserving perturbations can be defined directly on the simplex. In continuous state spaces, they can instead be constructed through transport maps. The remaining question is then one of representation: a probability measure must be described through a finite collection of relevant functionals, such as moments, interaction fields, or learned features. Since the perturbed population law is itself random, the score entering the gradient estimator is associated with the joint law of these functionals, rather than with the population measure directly. This viewpoint connects model-free MFC with representation learning on spaces of probability measures and provides the basis for treating discrete- and continuous-state problems within a common policy-gradient framework.


\section{Internship context}
This internship is carried out at the Centre de Mathématiques Appliquées (CMAP) of \'Ecole Polytechnique, in connection with the Institut Europlace de Finance. I am hosted within the probability group and the MathsFi team, whose research activities lie at the intersection of probability, stochastic control, mathematical finance, mean-field models, and numerical methods. The internship is supervised by Prof.~Huyên Pham. I also work closely with Samy Mekkaoui and Yadh Hafsi, with regular research meetings devoted to the progress and orientation of the project. On the academic side, the internship forms part of my final year at CentraleSupélec and of the MVA master's programme at \'Ecole Normale Supérieure Paris-Saclay.

The initial internship topic concerned the numerical resolution of mean-field games and mean-field control problems, potentially involving heterogeneous agents, using deep-learning methods. The subject was intentionally broad and was intended to leave room for the research direction to be refined during the internship. The first part of the work therefore consisted in exploring several approaches at the interface of mean-field control and machine learning before progressively identifying a more specific research question around reinforcement-learning methods for mean-field control. This evolution remained within the original scope of the internship while allowing the project to focus on a more clearly defined theoretical and numerical problem.

The project fits naturally within the broader research agenda of the host team. Mean-field control provides a framework for studying the optimal control of large populations of interacting agents and is closely connected to stochastic control, McKean--Vlasov dynamics, mathematical finance, and multi-agent systems. At the same time, reinforcement learning provides a complementary perspective in situations where the model is only partially known or where one wishes to construct algorithms directly from simulated or observed trajectories. The internship therefore lies at the intersection of several active research directions at CMAP: stochastic control, mean-field models, numerical probability, and machine learning.

This report reflects the state of the project at the time of the MVA presentation, in late August, while the internship itself continues until the end of October. Some theoretical developments and numerical conclusions should therefore be read as intermediate results, liable to be refined during the remaining part of the internship.

The internship is also closely connected with several courses from the MVA curriculum. The course \emph{Reinforcement Learning} by D.~Basu, E.~Kaufmann and O.~Maillard provides the most direct methodological background for the project, in particular for policy-gradient methods and model-free learning. \emph{Optimal Transport for Machine Learning} by G.~Peyré and J.~Delon is particularly relevant to the continuous-state extension, where perturbations of probability measures can naturally be formulated through transport maps. The course \emph{Stochastic Calculus in Machine Learning: Sampling and Generative Modeling} by A.~Oliviero-Durmus provides useful tools for working with stochastic processes, probability distributions and sampling-based methods, while \emph{Convex Optimization and Applications in Machine Learning} by A.~d'Aspremont provides a broader optimization perspective for the analysis and design of learning algorithms. The course \emph{Interactions} by J.~Randon-Furling is closely related to the modelling of systems of interacting agents and therefore to the mean-field perspective underlying the internship. Finally, \emph{Deep Learning} by V.~Lepetit and M.~Vakalopoulou provides the background for the use of neural-network parameterizations and function approximation in the numerical implementation of the proposed methods. Taken together, these courses reflect the interdisciplinary nature of the project, which combines reinforcement learning, probability, optimization, interacting-particle systems, optimal transport and deep learning.


\section{Objectives of the internship}

The central objective of the internship was to investigate whether policy-gradient methods for discrete-time mean-field control could be made fully model-free without relying on the finite-state logit representation of \cite{meunier2026model}. More precisely, the aim was to construct gradient estimators using only simulated trajectories, evaluations of the policy and observed costs, without differentiating or explicitly knowing the transition dynamics.

The first objective was theoretical. It consisted in identifying the additional dependence created by the population flow and in constructing a perturbation of this flow for which a likelihood-ratio representation could be derived. The perturbed gradient should approximate the gradient of the original MFC objective as the perturbation vanishes, while remaining estimable from samples. This also required studying the bias--variance trade-off introduced by the perturbation scale.

The second objective was to develop constructions adapted to both discrete and continuous state spaces. In the discrete case, the aim was to replace the logit perturbation by a probability-preserving perturbation defined directly on the simplex. In the continuous case, the objective was to perturb probability measures through transport maps and to express the required score through a finite-dimensional representation of the perturbed law. This raised a related representation-learning question: which functionals of the population distribution contain the information required by the dynamics, costs and policy, and how can the joint law of these functionals be estimated?

The third objective was algorithmic. The theoretical gradient representations had to be converted into practical Monte Carlo estimators, including an estimator of the sensitivity of the population flow with respect to the policy parameters. Particular attention was given to reusing simulated trajectories across the time horizon, controlling computational and memory costs, and selecting the perturbation scales. An adaptive procedure was also considered to reduce the dependence of the method on an environment-specific hyperparameter search.

Finally, the numerical objective was to evaluate the resulting methods on a range of mean-field control problems. The experiments were designed to compare the proposed estimators with standard REINFORCE and the finite-state MF-REINFORCE method under matched simulation budgets. They cover finite-state benchmarks, including cybersecurity, distribution planning and targeted advertising, as well as continuous-state problems involving linear--quadratic control, portfolio optimization and synchronization. Beyond final performance, the study examines gradient bias and variance, sensitivity to the perturbation scale, computational cost, and robustness to particle approximations of the population flow.


\section{Contributions}

The main contribution of this work is a perturbation-based framework for model-free policy gradients in discrete-time mean-field control. The central idea is to randomize the population argument so that its dependence on the policy parameters is transferred to the distribution of the perturbation. The gradient of the perturbed objective then decomposes into the usual policy-score term and an additional population-score term. Both terms admit likelihood-ratio representations and can therefore be estimated from simulated trajectories without differentiating the transition kernel.

The theoretical results accompanying this construction provide the main assurances needed for the method to be meaningful: the perturbed population laws remain controlled perturbations of the original flow, the perturbed objectives and gradients converge back to their unperturbed counterparts as the perturbation scale vanishes, and the Monte Carlo estimators are unbiased or consistent for the corresponding perturbed gradients. The bias--variance propositions further quantify the finite-sample trade-off created by the perturbation scale, clarifying why very small perturbations can be statistically unstable even though their asymptotic bias is small.

The contributions are organized as follows.

\paragraph{A transport view of perturbed population flows.}
We formulate perturbations directly at the level of probability measures. Given the population law $\mu_t^\theta=\P_{X_t^\theta}$, we introduce a random transport map $T_t^\lambda$ and replace the population argument by
\begin{align}
    M_t^{\lambda,\theta} = T_t^\lambda\sharp \mu_t^\theta.
\end{align}
The perturbation scale $\lambda$ controls the distance between the original and perturbed population laws. In continuous state spaces, this includes random maps of the form
\begin{align}
    T_t^\lambda(\omega,x)=x+\lambda f_t(\omega,x),
\end{align}
whereas in finite state spaces the same idea becomes a perturbation on the simplex. For both discrete and continuous MFC models, the estimator is driven by a random, probability-preserving deformation of the population law rather than by a model-specific coordinate system.

\paragraph{A likelihood-ratio correction for the mean-field dependence.}
The perturbed population argument is random, and its distribution depends on $\theta$ through $\mu_t^\theta$. The gradient of the perturbed objective therefore contains two terms: the usual REINFORCE policy-score term and an additional score term associated with the generated population argument. This second term is the correction missing from the classical REINFORCE formula in MFC. The report derives this decomposition and identifies the conditions under which it can be estimated without differentiating the transition kernel, the reward, or the population recursion.

\paragraph{A finite-dimensional representation for continuous laws.}
In a general continuous state space, there is no canonical dominating measure on $\Pc(\Xc)$, so the law of $M_t^{\lambda,\theta}$ need not have a usable density. We therefore introduce a representation map
\begin{align}
    \Psi(m)=\big(F_1(m),\ldots,F_d(m)\big)
\end{align}
and work with the random vector $Z_t^{\lambda,\theta}=\Psi(M_t^{\lambda,\theta})$. The population-score correction is then written through the density of $Z_t^{\lambda,\theta}$ in $\R^d$. This separates the policy-gradient construction from the choice of representation: in some examples $\Psi$ is given by exact sufficient statistics, such as moments, while in more general models it may consist of interaction fields, Fourier features, kernel features, or learned summaries of the population law.

\paragraph{A simplex perturbation for finite-state MFC.}
For $\mu\in\Delta_N$, we define the discrete transport perturbation
\begin{align}
    \Phi_\mu^\lambda(q)=(1-\lambda)\mu+\lambda q,
    \qquad q\in\Delta_N,
\end{align}
and use $M_t^{\lambda,\theta} = (1-\lambda)\mu_t^\theta+\lambda q_t$ as the perturbed population argument. In this setting, the simplex coordinates give an exact finite-dimensional representation. We derive the density and score of $M_t^{\lambda,\theta}$ explicitly, obtaining a model-free gradient estimator that acts directly on probability vectors. Compared with the logit perturbation of \cite{meunier2026model}, the construction does not require strictly positive population coordinates and is expressed in the geometry of the simplex itself. Under regularity assumptions, the perturbed gradient converges to the unperturbed gradient as $\lambda\to0$, with a first-order perturbation error.

\paragraph{A recursive sensitivity estimator and practical algorithms.}
The population-score term requires the sensitivity of the population flow with respect to the policy parameters. We derive a forward recursion for this quantity: the sensitivity at time $t$ depends only on sensitivities at earlier times, allowing the whole sensitivity flow to be estimated recursively from shared simulated trajectories. With cumulative computations, this gives a simulator cost that is linear in the horizon, in contrast with the quadratic cost of estimating each truncated sensitivity separately.

Finally, the estimators are implemented as fixed-scale and adaptive Monte Carlo algorithms. The adaptive version updates the main perturbation scale and the auxiliary sensitivity scale by comparing gradient estimates at nearby scales, aiming to balance perturbation bias and estimator variance without a full environment-specific grid search. The numerical study then evaluates these methods under matched simulator budgets on finite- and continuous-state benchmarks, comparing them with REINFORCE and MF-REINFORCE and reporting optimization performance, gradient diagnostics, runtime, and robustness to particle approximations of the population flow.


\section{Organization of the report}

Chapter~\ref{chapter:background} introduces the necessary background on Markov decision processes, policy-gradient methods, REINFORCE, and mean-field control, and positions the work relative to existing model-free MFC approaches. Chapter~\ref{chapter:methodology} develops the perturbation framework, the associated gradient representations, and the convergence and bias--variance results. Chapter~\ref{chapter:algorithms} turns these formulas into practical fixed-scale and adaptive Monte Carlo algorithms, including the recursive estimation of the population-flow sensitivity. Chapter~\ref{chapter:experiments} presents the numerical protocol and benchmark suite used to evaluate the methods. Chapter~\ref{chapter:discussion} discusses the main lessons, limitations, and open questions, and Chapter~\ref{chapter:conclusion} concludes the report.
